\paragraph{Theorem D.2.42 (Proved: the mixed six-, fifteen-, and thirty-torsion types).}
Consider the fourteen minimal patterns in (D.MT1), with the primitive
Mumford cap at every multiplicative site, the indicated additive
choice, and no additional regular logarithmic transformation.
Allow every finite multiplicity, affine twist and ordinary clutch
$e\gamma$ producing a simply connected manifold. Let
$M_{3,0}$ be the explicitly framed surgery of D.2.39, and
$Z_{5,\epsilon}$ the oriented type of D.2.36. Define
\[
\begin{aligned}
 Y_6&=M_{\mathrm{conic}}\# M_{3,0},\\
 Y_{15,\epsilon}&=M_{3,0}\# Z_{5,\epsilon},\\
 Y_{30}&=M_{\mathrm{conic}}\# M_{3,0}\# Z_{5,0}.
\end{aligned}
\]
The orientation-preserving diffeomorphism type is given by the
following table; zero positions have the source's zero-based indexing.
\[
\begin{array}{c|r|r|l|l|l}
\text{OS/profile}&a&b&\text{zero positions}&\text{additive choice}&\text{type}\\\hline
56&1&6&0&\text{none}&Y_6\\
56&5&6&3,4&\text{none}&Y_6\\
56/II&1&6&2&II\text{ original or log}&Y_6\\
56&1&15&2&\text{none}&Y_{15,0}\\
56&2&15&3,4&\text{none}&Y_{15,1}\\
56/II&1&15&0&II\text{ original or log}&Y_{15,0}\\
56/III&1&30&3&III\text{ log}&Y_{15,0}\\
56/III&2&15&3&III\text{ original}&Y_{15,1}\\
56&1&30&3,4&\text{none}&Y_{30}.
\end{array}                                               \tag{D.MT1}
\]
Thus the five $c=6$ patterns give one oriented type, the seven
$c=15$ patterns give two, and the two $c=30$ patterns give one.

\begin{proof}
\emph{The geometric input for every primary calculation.}
D.2.26 gives $H_2=H_3=\mathbb Z/c$ and $b_3=0$ in these
simply connected cases. As $H^2(X;\mathbb Z)=0$, each is spin.
There are twelve distinct marked initial rows. Write
$\Omega=a\,u^*w^*+b\,\gamma^*\delta^*$ in their actual markings.
The exact certificate below verifies its invariance and verifies
that both the multiplicative angular columns and the full images
of $A_i-I$ have integral Hermite lattice $\ker\gamma^*$.
Thus, after omitting the additive caps and one ordinary disc,
the common piece $W$ has torsion-free group
$\mathbb Z_\gamma\times F_r$. The primitive Mumford boundary
surjections and the independent remaining peripheral loops give
$H_2(V)\twoheadrightarrow H_2(W)$, with free first homology.
Consequently the degree-two comparisons on $W$ used in
D.2.35--36 apply to the actual marked classes here as well.

The three possible finite patterns have normal fibre coefficients
$6$, $15$, and $30/2=15$, respectively. All are divisible by
the stated residual $c$, as is the second normal coefficient,
a multiple of $b$. Every allowed original additive choice has
zero coefficient, and the clutch coefficient is $be$.
Hence the primitive fibre functional with zero swept values
descends modulo every prime dividing $c$, independently of
the allowed finite order and affine twist.

We record the local inputs to this comparison. The actual Tate
periods are $H=(n,ar),F=(-br,\rho)$; at the zero
$n\rho+ab r^2=1$. The doubled symmetric line (D.OP4) has
integral support
\[
          \psi=-a\rho x^2-2ab rxy+bn y^2+ax,
                                                        \tag{D.MT2}
\]
with the two exact period differences $-2ax$ and
$2by+b\rho-ab r$. For an odd prime $p\mid c$, all its
vertical slopes are divisible by $p$, and half its signed
wall degree is the unit $-a(\rho-br)$ modulo $p$.
At a rank-one site write $H=\alpha u_0,F=\beta u_0$, with
$u_0=(u_1,u_2)$ primitive and $\gcd(\alpha,\beta)=1$.
Its integral doubled support is
$-(a u_1/\alpha)j(j-\alpha)$; the extracted identities prove
$\alpha\mid a u_1$ and give both prescribed period differences.
Thus the auxiliary line extends on all the specified caps,
with the actual boundary marking.

For each $p\mid c$ there is exactly one additional
$C_p$-fixed elliptic double curve $E_p$, at an $I_p$ site.
Its primitive short period is $(p,ar)$ and its doubled support
is $-ap j(j-1)$. The line therefore has zero $C_p$ character
on both the rational fixed curve $C$ and $E_p$, for odd $p$.
The normal weights at $E_p$ are $w,-w$, where
$w=(ar)^{-1}\pmod p$. All these statements, including the
uniqueness of each fixed elliptic curve, are verified using
the actual primitive period columns in the twelve rows.
The finite caps have free scalar action by primitivity of
their gamma numerators; the original additive caps have
free scalar action. No extra prime-isotropy stratum is omitted.

\emph{The full three-primary class.}
For every row take $u_3\in H^2(X;\mathbb F_3)$ normalized
to have degree one on $C$. It lifts to $S^1$-equivariant
cohomology because $H^1(X;\mathbb F_3)=0$ and
$H^3(BS^1;\mathbb F_3)=0$. Normalize its point character
on $C$ to zero. The torsion-free group of $W$ and the
zero characters of (D.MT2) then make its $C_3$ point
character on $E_3$ zero as well, exactly as in D.2.39.
The restrictions of this lift are a class pulled back from
$C$ and an external product of degree-one classes on
$BC_3\times S^1\simeq(E_3)_{S^1}$.

Thomas's exact odd-prime external-product formula and
odd-degree vanishing make the Pontryagin third power of
both restrictions zero in coefficients $\mathbb Z/9$
\cite[equations (2.3), (2.5), p.267]{ffThomas56}.
The complementary locally free circle action has all
stabilizers of order prime to three. Its Borel pair has
the cohomology of a five-dimensional orbifold pair with
these coefficients, so its relative degree-six group
is zero. Thus restriction to $C\sqcup E_3$ is injective
in equivariant degree six, as proved in (D.T37)--(D.T38).
The full Pontryagin third power on $X$ is consequently zero.

The target calculation in part~1 of D.2.39 and Zhubr's
Theorem~5 now show that the entire spin-bordism image of
$X\to K(\mathbb Z/3,2)$ is zero. This use concerns an
arbitrary classifying map: it does not assume that
$H_2(X)$ itself equals $\mathbb Z/3$. It determines the
hidden subgroup of order three in $\mathbb Z/9$, in
addition to the ordinary mod-three cube. The primary
class is therefore exactly the class of $M_{3,0}$.

\emph{The two- and five-primary classes.}
When $c$ is even, $a$ is odd and $b/2$ is odd in every
row of (D.MT1). The unsquared line (D.C27) has the
monodromy-invariant semicharacter $a(u+w+uw)$.
Half of (D.MT2) and half of its rank-one support are
integral. Their scalar involution characters are one
on $C$ and zero on $E_2$, and their degree on $C$ is odd.
The common-piece comparison and the same simultaneous
Borel normalization as in D.2.35 thus give
$\langle u_2^3,[X]_2\rangle=1$. Its entire two-primary
bordism class is the conic-surgery class of D.2.8.

For $c=15$ or $30$, the five-primary fixed curves have
the zero characters supplied by (D.MT2). Normalize
$u_5[C]=1$. The actual localization of D.2.36 gives
\[
 \langle u_5^3,[X]\rangle=0,\qquad
 \mathfrak p_5=(w^2-1)/2\pmod5.                           \tag{D.MT3}
\]
The first equality detects the entire five-primary
fundamental class by D.1.10, and the second is the
complete refined Pontryagin class since four is a
unit modulo five. The actual weight is $w=4$ for
all $a=1$ rows, and $w=2$ for the $a=2$ rows.
Thus their complete bordism orbits are respectively
those of $Z_{5,0}$ and $Z_{5,1}$, by Zhubr's Theorem~5.

\emph{The complete oriented comparison.}
Iterate the prime splitting (D.MP9)--(D.MP10).
Positive integral homology of the distinct primary
Eilenberg--Mac Lane factors has zero cross tensor and
Tor products. The wedge-to-product inclusion is therefore
a simply connected homology equivalence and hence a
homotopy equivalence. It identifies the full spin-bordism
class over $K(\mathbb Z/c,2)$ with the primary classes
just computed, through the actual Chinese remainder
projections. The automorphism group splits in the same way.

The stated connected sums have exactly these classes:
their natural maps obtained by collapsing the separating
spheres induce the direct-sum isomorphism on $H_2$, and
the spin one-handle bordism identifies connected sum
with addition. They are simply connected and spin, have
$H_2=\mathbb Z/c$ and $b_3=0$. The uniqueness assertion
of Zhubr's Theorem~3 \cite[p.795]{manZhubr75} consequently
gives each diffeomorphism in (D.MT1). The two $c=15$
types are distinct because their five-primary refined
classes have different automorphism orbits.
\end{proof}

The extraction program was run at residuals $6,15,30$.
The independent exact local certificate is
\path{code/manifolds_three_mixed_AH.py}; its output records
all fourteen patterns, twelve marked rows, $464$
semicharacter substitutions and all three finite normal
coefficients. The full three-primary value is justified
by the primary operation and the Borel comparison above,
not by those finite algebraic checks alone.
